\clearpage
\SourceLeaf{279}
\renewcommand{\BookRunningTitle}{The Number of Classes and Their Determinant}
\renewcommand{\Article}[1]{\par\medskip\Needspace{4\baselineskip}\pdfbookmark[1]{Art. #1}{nexus-prior-art-#1}\centerline{\textbf{#1.}}\nobreak\smallskip}
\pdfbookmark[0]{The Number of Classes and Their Determinant. I}{nexus-prior}
\begin{center}
[I.]\par\medskip
\Large\bfseries ON THE CONNECTION BETWEEN THE NUMBER OF CLASSES INTO WHICH BINARY QUADRATIC FORMS ARE DISTRIBUTED AND THEIR DETERMINANT.\par\medskip
\normalsize FIRST PAPER\par\medskip
PRESENTED TO THE ROYAL SOCIETY 1834\ .\ .\ .\ .
\end{center}
\BookRule{.18\textwidth}

\Article{1}
Thirty-three years have now passed since we discovered the principles of the remarkable connection to which this paper is devoted, as was already announced at the end of the \emph{Arithmetical Investigations}. Other occupations, however, had long drawn us away from this inquiry, until more recently we were able to return to it and extend it through renewed efforts. Yet, since this new part of higher arithmetic exceeds the bounds of a single paper, this first paper will be devoted to forms with negative determinants; forms with positive determinants, which require an entirely distinct treatment, must remain reserved for another paper.

\Article{2}
The basis of the whole subject is a particular investigation of the number of all combinations of integer values taken by two indeterminate integers $x,y$ within a prescribed region. Evidently this problem can also be presented geometrically, by finding the number of \emph{complex numbers} whose representations fall within a prescribed figure. The nature of the figure will depend on that of the line surrounding it, and thus either on a single equation between the coordinates $x,y$ (when the boundary is a curve returning to itself) or on several such
\SourceLeaf{280}
equations (when it .consists of several curved or straight parts); it will be at our discretion whether to include in the count, or exclude from it, points corresponding to complex integers that may lie on the boundary itself.

In the analytic presentation of the problem, these conditions of restriction can always be expressed by requiring one or more given functions of the variables $x,y$, namely $P,Q,R$, etc., to take positive or nonnegative values (according as the value $0$ is excluded or admitted).

Thus, for example, if the prescribed figure is a circle with radius $=\surd A$, while its centre lies at a point corresponding to a complex integer, the analytic condition will be that $A-xx-yy$ is not negative, provided that, as we shall always suppose, we choose to retain points lying on the boundary itself. \mbox{If the figure is a triangle}, the three linear functions $ax+by+c$, $a'x+b'y+c'$, $a''x+b''y+c''$ must have nonnegative values, and similarly in other cases.

\Article{3}
Generally speaking, an \emph{exact} solution of the problem must proceed by first using the nature of the conditions to confine one variable, for example $x$, within limits between which the individual integer values are then taken successively. We must find how many integer values of the other variable $y$ correspond to each of them, and then add these numbers together. In special cases there will usually be particular devices for shortening the work.

For example, if the figure is, as above, a circle with radius $=\surd A$, let $r$ be the integer immediately below $\surd A$, or $\surd A$ itself if $A$ is a square. Likewise let $r',r'',r'''$, etc., $r^{(r)}$ be the integers immediately below $\surd(A-1)$, $\surd(A-4)$, $\surd(A-9)$, etc., as far as $\surd(A-rr)$. The required number will then be
\[
\begin{aligned}
&=2r+1+2(2r'+1)+2(2r''+1)+2(2r'''+1)+\text{etc.}\\
&=1+4r+4r'+4r''+4r'''+\text{etc.}+4r^{(r)}
\end{aligned}
\]

The following method is shorter in this example. Let $q$ be the integer immediately below $\surd\tfrac12 A$ (or equal to it when it is an integer), and let $r^{(q+1)}$, $r^{(q+2)}$, $r^{(q+3)}$, etc., be the integers immediately below $\surd(A-(q+1)^2)$, $\surd(A-(q+2)^2)$, $\surd(A-(q+3)^2)$, etc., as far as $\surd(A-rr)$. The required number will then be
\[
=4qq+1+4r+8(r^{(q+1)}+r^{(q+2)}+r^{(q+3)}+\text{etc.}+r^{(r)})
\]

\SourceLeaf{281}
This formula gave the following numbers:
\[
\begin{array}{r|r||r|r||r|r}
A&&A&&A&\\\hline
100&317&1000&3149&10000&31417\\
200&633&2000&6293&20000&62845\\
300&949&3000&9425&30000&94237\\
400&1257&4000&12581&40000&125629\\
500&1581&5000&15705&50000&157093\\
600&1885&6000&18853&60000&188453\\
700&2209&7000&21993&70000&219901\\
800&2521&8000&25137&80000&251305\\
900&2821&9000&28269&90000&282697\\
1000&3149&10000&31417&100000&314197
\end{array}
\]

\Article{4}
Our purpose does not require an exact determination, but rather the search for an expression that can approach the exact number as closely as desired as the limits are extended to infinity. But first of all, since this involves something vague, the matter must be stated more exactly.

We shall therefore suppose that the functions $P,Q,R$, etc., besides the variables $x,y$, involve a constant element $k$, so that each of $P,Q,R$, etc., is a homogeneous function of the three quantities $x,y,k$. In this way the figure determined by the equations $P=0,Q=0,R=0$, etc., will depend on $k$, so that different values of $k$ correspond to similar figures similarly situated with respect to the origin of coordinates; corresponding linear dimensions will be proportional to the values of $k$, and areas to the values of $kk$. Let the number of points within the figure now be denoted by $M$, and the area by $V$. It is evident that $M$ and $V$ must increase as $k$ increases; as $k$ increases to infinity, $M$ and $V$ will approach a ratio of equality as closely as desired. Or, if you require elementary clarity, given a quantity $\lambda$, however small, a threshold can always be assigned such that, for every value of $k$ exceeding it, $\frac MV$ must certainly lie between $1-\lambda$ and $1+\lambda$. In the customary manner this may be expressed by saying that $M=V$ for an infinite value of $k$.

In our example the required condition holds on putting $k=\surd A$, and the curve becomes a circle whose area $=\pi A$, where $\pi$ denotes the half-circumference of a circle of radius $=1$. The numbers given above clearly indicate the convergence.

\SourceLeaf{282}
Moreover, if it were worth the effort, we could easily complete the proof of this theorem with the rigour of the ancients; here, however, we have preferred to omit it, hastening on to more difficult matters.

\Article{5}
In this paper the boundary will be expressed by a \emph{single} equation of the form $axx+2bxy+cyy=A$, where $a,b,c$ are integers and $bb-ac$ is a negative number, which we shall put $=-D$. Evidently the curve defining the figure will be an ellipse, and it is readily seen that the squares of the semiaxes are the roots of the equation
\[
(ac-bb)qq-(a+c)Aq+AA=0\quad\text{or}\quad
=A\left(\frac{a+c\pm\surd(4bb+(a-c)^2)}{2(ac-bb)}\right)
\]

The product of these roots is $\frac{AA}{ac-bb}=\frac{AA}D$, and hence the area of the ellipse $=\frac{\pi A}{\surd D}$. We therefore infer that the number of all combinations of integer values of $x,y$ for which $axx+2bxy+cyy$ does not exceed the value $A$, as $A$ increases, comes ever closer to $\frac{\pi A}{\surd D}$, and for infinite $A$ must be set equal to that value. Moreover, it is evident that in this respect it makes no difference whether the combination $x=0,y=0$ is counted with the others or excluded. In this way the required number (in the latter sense) is simply the sum of the numbers of representations of the individual numbers $1,2,3\ldots A$ by the binary quadratic form $axx+2bxy+cyy$; and since some of those numbers cannot be represented by this form at all, while others admit more or fewer representations, the quantity $\frac\pi{\surd D}$ is to be regarded as the mean value of the number of representations of an indeterminate positive number by any form whose determinant $=-D$.
